\section{Discussion}
\label{sec:discussion}

\subsection{Accuracy of the Identification}
\label{sec:disc:accuracy}

The describing function approximation used in \cref{sec:theory:df} assumes
that the plant sufficiently attenuates higher harmonics of the relay switching
signal, so that the fundamental frequency component dominates.
This assumption holds well for plants with a roll-off of at least
\SI{-40}{\decibel\per decade} above $\omega_u$.
For thermal processes (which typically roll off much more steeply), the
approximation is highly accurate.

For plants with resonances near $\omega_u$, the describing function prediction
can be less accurate.
In such cases, it is recommended to add a low-pass filter in series with the
measurement before computing $\Au$ and $\Tu$.

\subsection{Robustness to Measurement Noise}
\label{sec:disc:noise}

Measurement noise directly affects the accuracy of $\Au$ and, more critically,
the switching instants used to compute $\Tu$.
A practical mitigation strategy is to:
\begin{enumerate}
    \item Use an adequate hysteresis $d \gg \sigma_n$, where $\sigma_n$ is
          the standard deviation of the measurement noise.
    \item Average over at least five oscillation cycles (as prescribed in
          \cref{sec:method:phase2}).
    \item Apply a moving-average pre-filter with window length
          $w = \Tu / 10$ before detecting zero-crossings.
\end{enumerate}

\subsection{Extension to SOPDT Plants}
\label{sec:disc:sopdt}

The FOPDT identification described in this paper is intentionally conservative:
it maps any higher-order plant onto a simple FOPDT model.
For plants where a second time constant is significant (such as
\eqref{eq:plant}), the estimated dead time $\hat{L}$ absorbs part of the
lag contributed by $T_2$, leading to a slightly conservative (i.e.\
detuned) PID.
This is generally desirable for robust industrial control.
An extension using a second-order relay feedback test~\cite{Wang1999}
is left for future work.

\subsection{Robustness to Mild Plant Nonlinearities}
\label{sec:disc:nonlinear}

The KZV derivation assumes a linear FOPDT plant, yet it retains a useful degree
of robustness to mild operating-point-dependent parameter variation.
Because the Bang-Bang limit cycle oscillates around a single setpoint~$r$
with amplitude~$\Au$, the identification is effectively performed on the
\emph{local linearisation} of the plant at~$r$: in the limit of small
hysteresis $d \ll \Au$, the estimates $(\hat{K}, \hat{T}, \hat{L})$ converge
to the true local FOPDT parameters at the current operating point.
For plants with moderate nonlinearities --- such as the gain and time-constant
variation typical of thermal and cooling circuits --- this means that repeating
the KZV identification at any new setpoint yields correctly updated PID
parameters without requiring an explicit nonlinear model.
A formal extension to gain-scheduled tuning across a range of operating
points is left for future work.

\subsection{Limitations}
\label{sec:disc:limitations}

\begin{itemize}
    \item The procedure requires the plant to be stable and to sustain a
          limit cycle under Bang-Bang control.
          Integrating or unstable plants require a modified relay feedback
          scheme.
    \item Nonlinear plant behaviour (e.g.\ hysteresis, saturation) may cause
          the limit cycle to be asymmetric, leading to biased estimates.
    \item The method is derived for SISO systems. Multi-variable extensions
          are not considered here.
\end{itemize}

\subsection{Comparison with Related Methods}
\label{sec:disc:related}

The KZV is most closely related to the relay feedback tuning method of
\AA{}str\"om and H\"agglund~\cite{Astrom1984}, which also uses a relay
(Bang-Bang without hysteresis) to excite a sustained oscillation.
The two methods share the same underlying experiment but make different
trade-offs:

\begin{itemize}
    \item \textbf{\AA{}str\"om \& H\"agglund} read a single point on the
          Nyquist curve ($\Ku$, $\omega_u$) from the limit cycle via
          describing-function analysis. Because this requires essentially no
          assumption about the plant structure, the method is broadly
          applicable across very different process types. PID parameters are
          then obtained via Ziegler--Nichols-style rules, which do not
          explicitly account for dead time.
    \item \textbf{KZV} extracts three time-domain slope measurements from
          the same limit cycle to estimate the full FOPDT triple
          $(\hat{K}, \hat{T}, \hat{L})$ and tunes via IMC, which exposes a
          robustness knob $\lambda$ and treats dead time as a first-class
          parameter. The price for this richer model is two-fold: (i) the
          plant must be reasonably well approximated by FOPDT, and (ii)
          slope estimates are more sensitive to measurement noise than
          amplitude/period estimates.
\end{itemize}

The two approaches are therefore best viewed as \emph{complementary}.
For general-purpose auto-tuning across heterogeneous plants, classical
relay feedback remains the more conservative default. For the FOPDT-friendly
process class targeted in this work --- in particular thermal and cooling
systems with non-negligible dead time --- the KZV's explicit dead-time
estimate and IMC tuning typically yield a more transparent and easier-to-tune
result from the same experimental data.
